% ================== UNIT V ==================
\unitbanner{V}{Learning, Expert Systems and Applications of AI}{06 Hours \textbullet\ CO204.5}

\begin{outcome}
\begin{itemize}
\item distinguish supervised, unsupervised and reinforcement learning \btl{L2}
\item build a decision tree using entropy and information gain \btl{L3}
\item explain artificial neurons and train a perceptron by hand \btl{L3}
\item describe the architecture of an expert system and design a rule base \btl{L2/L6}
\item illustrate AI applications in NLP, vision and robotics and evaluate ethical issues \btl{L4/L5}
\end{itemize}
\end{outcome}

\section{Introduction to Learning}
An agent is \textbf{learning} if it improves its performance on future tasks after making observations about the world (Mitchell: \emph{a program learns from experience E w.r.t.\ task T and performance measure P if its performance at T, as measured by P, improves with E}).

\subsection{Forms of Learning}
\begin{center}
\begin{tikzpicture}[r/.style={draw=navy,thick,rounded corners,fill=#1,text width=4.4cm,minimum height=3.3cm,align=left,font=\small}]
\node[r=goldL](s){\textbf{\color{navy}Supervised}\\ Labelled pairs $(x,y)$; learn $h\approx f$.\\ \emph{Classification} (discrete $y$), \emph{regression} (continuous $y$).\\ e.g.\ spam filter, house price.};
\node[r=tealL,right=0.25cm of s](u){\textbf{\color{navy}Unsupervised}\\ Unlabelled $x$; discover structure.\\ \emph{Clustering}, dimensionality reduction, association rules.\\ e.g.\ customer segmentation.};
\node[r=leafL,right=0.25cm of u](r){\textbf{\color{navy}Reinforcement}\\ Agent acts, receives \emph{rewards/punishments}; learns a policy $\pi$ maximising return.\\ e.g.\ AlphaGo, robot walking.};
\end{tikzpicture}
\end{center}
Also: \emph{semi-supervised}, \emph{self-supervised} (foundation models) and \emph{transfer} learning. Key concepts: hypothesis space, \textbf{generalisation}, \textbf{overfitting} vs.\ \textbf{underfitting}, training/validation/test split, cross-validation, bias--variance trade-off, Ockham's razor.

\section{Learning by Decision Trees}
A decision tree represents a function mapping attribute vectors to a decision; each internal node tests an attribute, branches are attribute values, leaves are outputs. Learned top-down greedily (ID3/C4.5/CART), choosing the \emph{most important} attribute first.
\subsection{Entropy and Information Gain}
\[H(S)=-\sum_{k}p_k\log_2p_k,\qquad \textit{Gain}(S,A)=H(S)-\sum_{v\in\textit{Values}(A)}\frac{|S_v|}{|S|}H(S_v)\]
\begin{example}[Play Tennis (14 examples: 9 Yes, 5 No)]
\small
$H(S)=-\tfrac{9}{14}\log_2\tfrac{9}{14}-\tfrac{5}{14}\log_2\tfrac{5}{14}=0.940$.\\
\textbf{Outlook}: Sunny (2Y,3N) $H=0.971$; Overcast (4Y,0N) $H=0$; Rain (3Y,2N) $H=0.971$.\\
$\textit{Gain}(S,\textit{Outlook})=0.940-\tfrac{5}{14}(0.971)-\tfrac{4}{14}(0)-\tfrac{5}{14}(0.971)=\mathbf{0.247}$.\\
Similarly Gain(Humidity)$=0.151$, Gain(Wind)$=0.048$, Gain(Temp)$=0.029$ $\Rightarrow$ \textbf{Outlook} is the root.\\[2pt]
\begin{tikzpicture}[level distance=1.15cm,level 1/.style={sibling distance=3.4cm},level 2/.style={sibling distance=1.7cm},
 every node/.style={draw=navy,rounded corners,fill=goldL,font=\scriptsize\bfseries},edge from parent/.style={draw=navy,thick},
 yes/.style={fill=leaf,text=white},no/.style={fill=crimson,text=white}]
\node[fill=navy,text=white]{Outlook}
 child{node{Humidity} child{node[no]{No} edge from parent node[left,draw=none,fill=none]{High}} child{node[yes]{Yes} edge from parent node[right,draw=none,fill=none]{Normal}} edge from parent node[left,draw=none,fill=none]{Sunny}}
 child{node[yes]{Yes} edge from parent node[right,draw=none,fill=none]{Overcast}}
 child{node{Wind} child{node[no]{No} edge from parent node[left,draw=none,fill=none]{Strong}} child{node[yes]{Yes} edge from parent node[right,draw=none,fill=none]{Weak}} edge from parent node[right,draw=none,fill=none]{Rain}};
\end{tikzpicture}\\
Rule extraction: \texttt{IF Outlook=Sunny AND Humidity=Normal THEN Yes}, \dots
\end{example}
Other measures: \textbf{Gini index} $1-\sum p_k^2$ (CART), \textbf{gain ratio} (C4.5). Overfitting is controlled by \emph{pre-pruning} (max depth, min samples) and \emph{post-pruning} ($\chi^2$ pruning, reduced-error pruning). Trees are interpretable --- a key property for explainable AI.

\section{Artificial Neural Networks and the Perceptron}
\subsection{The Artificial Neuron}
\begin{center}
\begin{tikzpicture}[i/.style={circle,draw=navy,fill=goldL,minimum size=0.75cm,font=\small},>=Latex]
\foreach \k/\y in {1/1.6,2/0.5,n/-1.1}{\node[i](x\k) at (0,\y){$x_{\k}$};}
\node at (0,-0.25){$\vdots$}; \node[i,fill=tealL](b) at (0,2.7){$1$};
\node[circle,draw=navy,very thick,fill=navy,text=white,minimum size=1.4cm](s) at (3.5,0.6){$\Sigma$};
\node[draw=navy,thick,fill=gold,minimum size=1cm,rounded corners](f) at (5.7,0.6){$g(\cdot)$};
\draw[->,thick](b)--node[above,font=\scriptsize]{$w_0$ (bias)}(s);
\foreach \k in {1,2,n}{\draw[->,thick](x\k)--node[above,font=\scriptsize,pos=0.4]{$w_{\k}$}(s);}
\draw[->,thick](s)--node[above,font=\scriptsize]{$in$}(f); \draw[->,very thick,crimson](f)--++(1.6,0) node[right]{$y$};
\end{tikzpicture}
\end{center}
$y=g\bigl(\sum_{i=0}^{n}w_ix_i\bigr)$ with activation $g$: \emph{step/threshold}, \emph{sigmoid} $1/(1+e^{-z})$, \emph{tanh}, \emph{ReLU} $\max(0,z)$. Networks: \emph{feed-forward} (multilayer perceptron) and \emph{recurrent}; trained by \textbf{back-propagation} with gradient descent. Deep learning stacks many layers (CNNs for vision, Transformers for language).

\subsection{The Perceptron (Rosenblatt, 1958)}
A single neuron with a hard threshold: $h_w(x)=1$ if $w\cdot x\ge0$ else $0$.
\begin{algo}[Perceptron learning rule]
Initialise $w$ randomly. For each training example $(x,y)$: \quad $w_i\leftarrow w_i+\eta\,(y-h_w(x))\,x_i$ \quad (learning rate $\eta$). Repeat until no errors.
\end{algo}
\textbf{Perceptron convergence theorem}: converges if the data are \emph{linearly separable}. It learns AND, OR, NOT but \textbf{not XOR} (Minsky \& Papert, 1969) --- solved by multilayer networks with a hidden layer.
\begin{example}[Training a perceptron for AND ($\eta=1$, $w=(w_0,w_1,w_2)=(0,0,0)$, $x_0=1$)]
\small
\begin{tabular}{cccccc}
\rowcolor{navy}\color{gold}$x_1$&\color{gold}$x_2$&\color{gold}$y$&\color{gold}$\hat y$&\color{gold}error&\color{gold}new $w$\\
0&0&0&1&$-1$&$(-1,0,0)$\\ \rowcolor{ivory}0&1&0&0&0&$(-1,0,0)$\\ 1&0&0&0&0&$(-1,0,0)$\\ \rowcolor{ivory}1&1&1&0&$+1$&$(0,1,1)$\\
\multicolumn{6}{l}{\dots\ after a few epochs: $w=(-2,1,1)$ i.e.\ $x_1+x_2\ge2$ --- correct AND boundary.}\\
\end{tabular}
\end{example}

\section{Expert Systems}
An expert system is a program that emulates the decision-making ability of a human expert in a narrow domain using a knowledge base of facts and heuristic rules. Classic systems: \textbf{MYCIN} (bacterial infections, certainty factors), \textbf{DENDRAL} (mass spectrometry), \textbf{XCON/R1} (configuring VAX computers), \textbf{PROSPECTOR} (mineral exploration).
\subsection{Architecture}
\begin{center}
\begin{tikzpicture}[c/.style={draw=navy,thick,rounded corners,fill=#1,minimum width=3.1cm,minimum height=1cm,align=center,font=\small},>=Latex]
\node[c=goldL](ui){User interface}; \node[c=navy,text=white,right=1.4cm of ui](ie){Inference engine};
\node[c=tealL,above right=0.3cm and 1.2cm of ie](kb){Knowledge base\\(rules + facts)};
\node[c=tealL,below right=0.3cm and 1.2cm of ie](wm){Working memory\\(case facts)};
\node[c=leafL,below=0.8cm of ui](ex){Explanation facility}; \node[c=leafL,above=0.8cm of ui](ka){Knowledge acquisition};
\node[left=0.6cm of ui,font=\small\bfseries](u){User}; \node[left=0.6cm of ka,font=\small\bfseries](ke){Expert / KE};
\draw[<->,thick](u)--(ui); \draw[<->,thick](ui)--(ie); \draw[<->,thick](ie)--(kb); \draw[<->,thick](ie)--(wm);
\draw[<->,thick](ui)--(ex); \draw[->,thick](ke)--(ka); \draw[->,thick,gold!80!black](ka)-|(kb); \draw[->,thick](ie)|-(ex);
\end{tikzpicture}
\end{center}
\subsection{Knowledge Acquisition}
The process of extracting, structuring and organising knowledge from domain experts, books and data --- performed by a \textbf{knowledge engineer} through interviews, protocol analysis, observation, repertory grids and machine learning. It is the main \textbf{bottleneck} in building expert systems.
\subsection{Inference Engine and Rule-Based Reasoning}
Knowledge is encoded as production rules: \texttt{IF <conditions> THEN <conclusion/action> [CF]}. The inference engine runs a \textbf{recognise--act cycle}: (1) \emph{match} rules against working memory (Rete algorithm), (2) \emph{conflict resolution} --- choose one rule by specificity, recency, priority or refractoriness, (3) \emph{act} --- fire it, updating working memory.
\begin{itemize}
\item \textbf{Forward chaining} (data-driven): monitoring, configuration, planning (XCON).
\item \textbf{Backward chaining} (goal-driven): diagnosis (MYCIN).
\item \textbf{Uncertainty}: certainty factors $CF\in[-1,1]$; $CF(H,E_1\wedge E_2)$ combined as $CF_1+CF_2(1-CF_1)$ for positive CFs.
\end{itemize}
\begin{example}[Mini medical rule base]
\small R1: IF fever AND cough THEN flu-suspected (CF 0.6)\\
R2: IF flu-suspected AND body-ache THEN influenza (CF 0.8)\\
R3: IF fever AND rash THEN measles-suspected (CF 0.7)\\
Facts: fever, cough, body-ache $\Rightarrow$ R1 fires $\Rightarrow$ R2 fires $\Rightarrow$ \textbf{influenza} with $CF=0.6\times0.8=0.48$. Explanation: ``\emph{because} fever and cough (R1), \emph{and} body-ache (R2)''.
\end{example}
Expert-system shells: CLIPS, Jess, Drools, Prolog-based shells. Advantages: consistency, availability, explanation; limitations: brittleness, no common sense, costly maintenance.

\section{Applications of AI}
\subsection{Natural Language Processing}
Levels: morphology $\to$ syntax (parsing) $\to$ semantics $\to$ pragmatics $\to$ discourse. Tasks: tokenisation, POS tagging, named-entity recognition, sentiment analysis, machine translation, question answering, summarisation, chatbots. Evolution: rule-based $\to$ statistical (n-grams, HMM) $\to$ word embeddings (word2vec) $\to$ Transformers / LLMs (BERT, GPT, Claude).
\subsection{Computer Vision}
Image formation $\to$ low-level processing (edges, filters) $\to$ features $\to$ recognition. Tasks: image classification, object detection (YOLO), semantic segmentation, face recognition, OCR, medical image analysis, autonomous driving perception. Deep CNNs and Vision Transformers dominate.
\subsection{Robotics}
Robots perceive (cameras, LiDAR, IMU), plan (path/motion planning, SLAM --- simultaneous localisation and mapping) and act (manipulators, mobile bases). Types: industrial arms, mobile robots, drones, humanoids, surgical robots (da Vinci), Mars rovers. Challenges: uncertainty, real-time control, human--robot interaction.

\section{Ethical and Societal Issues in AI}
\begin{center}\small
\begin{tabularx}{\linewidth}{>{\leavevmode\bfseries\color{navy}}p{3.4cm}X}
\toprule
Bias \& fairness & Models inherit societal bias from data (e.g.\ hiring tools, face recognition error rates across groups).\\
Privacy \& surveillance & Mass data collection, face recognition in public spaces; India's DPDP Act 2023, EU GDPR.\\
Transparency \& explainability & ``Black-box'' decisions in credit, health, justice; need for XAI (LIME, SHAP).\\
Accountability \& liability & Who is responsible when an autonomous car or medical AI errs?\\
Safety \& robustness & Adversarial examples, hallucinations, misuse; alignment of goals with human values.\\
Employment & Automation displaces routine jobs; reskilling and new job creation.\\
Misinformation & Deepfakes, synthetic media, election manipulation.\\
Autonomous weapons & Lethal autonomous systems and meaningful human control.\\
Environment & Energy and water footprint of training large models.\\\bottomrule
\end{tabularx}\end{center}
\textbf{Frameworks:} EU AI Act (risk-based regulation, 2024), OECD AI Principles, UNESCO Recommendation on AI Ethics, NITI Aayog \emph{Responsible AI for All} (India), IEEE Ethically Aligned Design. Core principles: \emph{beneficence, non-maleficence, autonomy, justice, explicability}.
\begin{keypoint}[Responsible AI checklist for every project]
Data consent \& provenance $\cdot$ bias audit across groups $\cdot$ explanation for affected users $\cdot$ human-in-the-loop for high-stakes decisions $\cdot$ security testing $\cdot$ documented limitations (model cards).
\end{keypoint}

\subsection*{Unit V --- Summary}
\begin{itemize}
\item Supervised, unsupervised and reinforcement learning differ in the feedback available.
\item Decision trees choose attributes by information gain; perceptrons learn linear boundaries.
\item Expert systems = knowledge base + inference engine + explanation + knowledge acquisition.
\item AI powers NLP, vision and robotics --- and must be developed responsibly.
\end{itemize}

% ---- Industry readiness (generated by 07_Interview_Prep/src/build_interview.py)
\section{Industry Readiness: Skills \& Interview Corner}
\begin{tcolorbox}[enhanced,colback=navy,colframe=gold,boxrule=1.2pt,arc=2mm,coltext=white,
 title={\textbf{Top industry skills you gain in Unit V}},fonttitle=\color{navy},colbacktitle=gold]
\begin{itemize}\setlength\itemsep{1pt}
\item Training, evaluating and interpreting decision-tree models
\item Explaining overfitting, bias--variance and validation strategy
\item Building a perceptron / small neural network and debugging training
\item Designing rule-based expert systems and chatbots with fallbacks
\item Auditing models for bias, privacy and explainability
\item Mapping NLP, vision and robotics problems to the right technique
\end{itemize}
\end{tcolorbox}
\begin{tabularx}{\linewidth}{>{\bfseries\color{navy}}p{3.4cm}X}
Tools to master & scikit-learn, PyTorch / Keras, SHAP / Fairlearn, Rasa or simple intent classifiers\\
Job roles & ML Engineer, Data Scientist, Conversational AI Developer, Responsible-AI Analyst\\
Portfolio output & An end-to-end classifier with a model card (data, metrics, fairness checks, limitations) and a deployed demo.\\
\end{tabularx}
\subsection*{Interview corner --- top AI company style questions}
{\small\color{slate}Practice questions modelled on common AI/ML interview formats; full set of 15 with model answers in \texttt{07\_Interview\_Prep}.}
\begin{tcolorbox}[enhanced,breakable,colback=white,colframe=teal,arc=2mm,boxrule=0.8pt,left=3mm,right=3mm,
 title={\textbf{QV.01}\quad Explain supervised, unsupervised and reinforcement learning with one industry example each.},fonttitle=\color{white}\small,colbacktitle=teal]
\colorbox{teal!15}{\scriptsize\bfseries\color{teal} CONCEPT}\
\colorbox{leaf!15}{\scriptsize\bfseries\color{leaf} EASY}\
\colorbox{navy!8}{\scriptsize\bfseries\color{navy} IT Services / GCC}\par\smallskip
\textbf{\color{navy}Model answer.} Supervised: credit-risk scoring from labelled defaults. Unsupervised: customer segmentation by clustering. Reinforcement: dynamic pricing or robot control learned from rewards.
\end{tcolorbox}
\begin{tcolorbox}[enhanced,breakable,colback=white,colframe=teal,arc=2mm,boxrule=0.8pt,left=3mm,right=3mm,
 title={\textbf{QV.02}\quad How does a decision tree choose a split? Compare entropy and Gini.},fonttitle=\color{white}\small,colbacktitle=teal]
\colorbox{teal!15}{\scriptsize\bfseries\color{teal} CONCEPT}\
\colorbox{gold!85!black!15}{\scriptsize\bfseries\color{gold!85!black} MEDIUM}\
\colorbox{navy!8}{\scriptsize\bfseries\color{navy} Product / Big Tech}\par\smallskip
\textbf{\color{navy}Model answer.} It picks the split with the largest impurity reduction. Entropy uses --$\Sigma$p log p; Gini uses 1 -- $\Sigma$p$^2$. They usually give similar trees; Gini is slightly cheaper to compute.
\end{tcolorbox}
\begin{tcolorbox}[enhanced,breakable,colback=white,colframe=teal,arc=2mm,boxrule=0.8pt,left=3mm,right=3mm,
 title={\textbf{QV.03}\quad Why can't a single perceptron learn XOR?},fonttitle=\color{white}\small,colbacktitle=teal]
\colorbox{teal!15}{\scriptsize\bfseries\color{teal} CONCEPT}\
\colorbox{gold!85!black!15}{\scriptsize\bfseries\color{gold!85!black} MEDIUM}\
\colorbox{navy!8}{\scriptsize\bfseries\color{navy} AI Research Lab}\par\smallskip
\textbf{\color{navy}Model answer.} XOR is not linearly separable; a perceptron draws one line. A hidden layer creates a non-linear boundary, so a 2-layer MLP solves it.
\end{tcolorbox}
\begin{tcolorbox}[enhanced,breakable,colback=white,colframe=crimson,arc=2mm,boxrule=0.8pt,left=3mm,right=3mm,
 title={\textbf{QV.10}\quad Design a content-moderation system for a social platform in Indian languages.},fonttitle=\color{white}\small,colbacktitle=crimson]
\colorbox{crimson!15}{\scriptsize\bfseries\color{crimson} DESIGN}\
\colorbox{crimson!15}{\scriptsize\bfseries\color{crimson} HARD}\
\colorbox{navy!8}{\scriptsize\bfseries\color{navy} Product / Big Tech}\par\smallskip
\textbf{\color{navy}Model answer.} Multilingual text and image classifiers, policy taxonomy, thresholds per category, human review for borderline cases, appeal flow, feedback loop for retraining, and metrics per language to catch bias.
\end{tcolorbox}
\begin{tcolorbox}[enhanced,breakable,colback=white,colframe=crimson,arc=2mm,boxrule=0.8pt,left=3mm,right=3mm,
 title={\textbf{QV.11}\quad Design an expert-system-plus-ML assistant for rural primary-health workers.},fonttitle=\color{white}\small,colbacktitle=crimson]
\colorbox{crimson!15}{\scriptsize\bfseries\color{crimson} DESIGN}\
\colorbox{crimson!15}{\scriptsize\bfseries\color{crimson} HARD}\
\colorbox{navy!8}{\scriptsize\bfseries\color{navy} HealthTech}\par\smallskip
\textbf{\color{navy}Model answer.} Rule-based triage protocol for safety-critical decisions, an ML model for risk ranking, offline-first mobile app, local-language voice UI, escalation to doctors, and audit logs; validate with clinicians.
\end{tcolorbox}

% ---- Unit quiz (generated by 11_Quizzes/src/build_quizzes.py)
\section{Unit Quiz --- test yourself}
{\small\color{slate}Fifteen multi-step questions. Answer before checking the key on the next page. Self-check notebook: \texttt{11\_Quizzes/Quiz\_SelfCheck.ipynb}.}\par\medskip
\begin{tcolorbox}[enhanced,breakable,colback=white,colframe=navyL,arc=2mm,boxrule=0.7pt,left=3mm,right=3mm,top=2mm,bottom=2mm]
\textbf{\color{navy}QV.01}\quad A dataset has 9 positive and 5 negative examples. What is its entropy?\hfill\btl{L3}\par\smallskip
{\small \textbf{\color{navyL}(A)}\ 0.500\\[1pt]\textbf{\color{navyL}(B)}\ 0.643\\[1pt]\textbf{\color{navyL}(C)}\ 0.940\\[1pt]\textbf{\color{navyL}(D)}\ 1.000}
\end{tcolorbox}
\begin{tcolorbox}[enhanced,breakable,colback=white,colframe=navyL,arc=2mm,boxrule=0.7pt,left=3mm,right=3mm,top=2mm,bottom=2mm]
\textbf{\color{navy}QV.02}\quad Outlook splits Play Tennis into Sunny (2+, 3--), Overcast (4+, 0--) and Rain (3+, 2--). What is Gain(S, Outlook)?\hfill\btl{L3}\par\smallskip
{\small \textbf{\color{navyL}(A)}\ 0.048\\[1pt]\textbf{\color{navyL}(B)}\ 0.151\\[1pt]\textbf{\color{navyL}(C)}\ 0.247\\[1pt]\textbf{\color{navyL}(D)}\ 0.940}
\end{tcolorbox}
\begin{tcolorbox}[enhanced,breakable,colback=white,colframe=navyL,arc=2mm,boxrule=0.7pt,left=3mm,right=3mm,top=2mm,bottom=2mm]
\textbf{\color{navy}QV.03}\quad What is the Gini impurity of a node with 9 positive and 5 negative examples?\hfill\btl{L3}\par\smallskip
{\small \textbf{\color{navyL}(A)}\ 0.230\\[1pt]\textbf{\color{navyL}(B)}\ 0.459\\[1pt]\textbf{\color{navyL}(C)}\ 0.541\\[1pt]\textbf{\color{navyL}(D)}\ 0.940}
\end{tcolorbox}
\begin{tcolorbox}[enhanced,breakable,colback=white,colframe=navyL,arc=2mm,boxrule=0.7pt,left=3mm,right=3mm,top=2mm,bottom=2mm]
\textbf{\color{navy}QV.04}\quad A perceptron has w = (w$_0$, w$_1$, w$_2$) = (--1, 0, 0) with bias input x$_0$ = 1. On x = (1, 1) with target y = 1 and $\eta$ = 1, what are the weights after one update?\hfill\btl{L3}\par\smallskip
{\small \textbf{\color{navyL}(A)}\ (--1, 0, 0)\\[1pt]\textbf{\color{navyL}(B)}\ (0, 1, 1)\\[1pt]\textbf{\color{navyL}(C)}\ (--2, --1, --1)\\[1pt]\textbf{\color{navyL}(D)}\ (1, 1, 1)}
\end{tcolorbox}
\begin{tcolorbox}[enhanced,breakable,colback=white,colframe=navyL,arc=2mm,boxrule=0.7pt,left=3mm,right=3mm,top=2mm,bottom=2mm]
\textbf{\color{navy}QV.05}\quad What is the smallest feed-forward network with threshold units that represents XOR?\hfill\btl{L4}\par\smallskip
{\small \textbf{\color{navyL}(A)}\ One perceptron\\[1pt]\textbf{\color{navyL}(B)}\ One hidden layer with 2 units\\[1pt]\textbf{\color{navyL}(C)}\ One hidden layer with 4 units\\[1pt]\textbf{\color{navyL}(D)}\ Two hidden layers with 1 unit each}
\end{tcolorbox}
\begin{tcolorbox}[enhanced,breakable,colback=white,colframe=navyL,arc=2mm,boxrule=0.7pt,left=3mm,right=3mm,top=2mm,bottom=2mm]
\textbf{\color{navy}QV.06}\quad A classifier has TP = 40, FP = 10 and FN = 20. What are its precision, recall and F1?\hfill\btl{L3}\par\smallskip
{\small \textbf{\color{navyL}(A)}\ 0.80, 0.67, 0.73\\[1pt]\textbf{\color{navyL}(B)}\ 0.67, 0.80, 0.73\\[1pt]\textbf{\color{navyL}(C)}\ 0.80, 0.80, 0.80\\[1pt]\textbf{\color{navyL}(D)}\ 0.57, 0.67, 0.62}
\end{tcolorbox}
\begin{tcolorbox}[enhanced,breakable,colback=white,colframe=navyL,arc=2mm,boxrule=0.7pt,left=3mm,right=3mm,top=2mm,bottom=2mm]
\textbf{\color{navy}QV.07}\quad A decision tree scores 99\% training accuracy and 70\% test accuracy; a depth-3 tree scores 85\% and 84\%. What is the right reading?\hfill\btl{L4}\par\smallskip
{\small \textbf{\color{navyL}(A)}\ The deep tree is better\\[1pt]\textbf{\color{navyL}(B)}\ The deep tree overfits; prefer the shallow tree or prune\\[1pt]\textbf{\color{navyL}(C)}\ Both underfit\\[1pt]\textbf{\color{navyL}(D)}\ The test set is wrong}
\end{tcolorbox}
\begin{tcolorbox}[enhanced,breakable,colback=white,colframe=navyL,arc=2mm,boxrule=0.7pt,left=3mm,right=3mm,top=2mm,bottom=2mm]
\textbf{\color{navy}QV.08}\quad Two independent rules support the same conclusion with CF$_1$ = 0.6 and CF$_2$ = 0.5. What is the combined CF (MYCIN)?\hfill\btl{L3}\par\smallskip
{\small \textbf{\color{navyL}(A)}\ 0.30\\[1pt]\textbf{\color{navyL}(B)}\ 0.55\\[1pt]\textbf{\color{navyL}(C)}\ 0.80\\[1pt]\textbf{\color{navyL}(D)}\ 1.10}
\end{tcolorbox}
\begin{tcolorbox}[enhanced,breakable,colback=white,colframe=navyL,arc=2mm,boxrule=0.7pt,left=3mm,right=3mm,top=2mm,bottom=2mm]
\textbf{\color{navy}QV.09}\quad Rule: IF fever (CF 0.9) AND cough (CF 0.7) THEN flu with rule CF 0.8. What is CF(flu)?\hfill\btl{L3}\par\smallskip
{\small \textbf{\color{navyL}(A)}\ 0.50\\[1pt]\textbf{\color{navyL}(B)}\ 0.56\\[1pt]\textbf{\color{navyL}(C)}\ 0.63\\[1pt]\textbf{\color{navyL}(D)}\ 0.72}
\end{tcolorbox}
\begin{tcolorbox}[enhanced,breakable,colback=white,colframe=navyL,arc=2mm,boxrule=0.7pt,left=3mm,right=3mm,top=2mm,bottom=2mm]
\textbf{\color{navy}QV.10}\quad Which pairing of classic expert system and chaining strategy is correct?\hfill\btl{L4}\par\smallskip
{\small \textbf{\color{navyL}(A)}\ MYCIN: forward; XCON: backward\\[1pt]\textbf{\color{navyL}(B)}\ MYCIN: backward; XCON: forward\\[1pt]\textbf{\color{navyL}(C)}\ Both backward\\[1pt]\textbf{\color{navyL}(D)}\ Both forward}
\end{tcolorbox}
\begin{tcolorbox}[enhanced,breakable,colback=white,colframe=navyL,arc=2mm,boxrule=0.7pt,left=3mm,right=3mm,top=2mm,bottom=2mm]
\textbf{\color{navy}QV.11}\quad In a dataset with 2\% positives, a model predicts 'negative' for everyone. What is its accuracy, and what should you report instead?\hfill\btl{L5}\par\smallskip
{\small \textbf{\color{navyL}(A)}\ 2\%; report accuracy\\[1pt]\textbf{\color{navyL}(B)}\ 98\%; report precision, recall, F1 and PR-AUC\\[1pt]\textbf{\color{navyL}(C)}\ 50\%; report ROC-AUC only\\[1pt]\textbf{\color{navyL}(D)}\ 100\%; nothing else needed}
\end{tcolorbox}
\begin{tcolorbox}[enhanced,breakable,colback=white,colframe=navyL,arc=2mm,boxrule=0.7pt,left=3mm,right=3mm,top=2mm,bottom=2mm]
\textbf{\color{navy}QV.12}\quad What are sigmoid(0), tanh(0) and ReLU(--2)?\hfill\btl{L3}\par\smallskip
{\small \textbf{\color{navyL}(A)}\ 0.5, 0, 0\\[1pt]\textbf{\color{navyL}(B)}\ 0, 0, --2\\[1pt]\textbf{\color{navyL}(C)}\ 1, 1, 0\\[1pt]\textbf{\color{navyL}(D)}\ 0.5, 0.5, 2}
\end{tcolorbox}
\begin{tcolorbox}[enhanced,breakable,colback=white,colframe=navyL,arc=2mm,boxrule=0.7pt,left=3mm,right=3mm,top=2mm,bottom=2mm]
\textbf{\color{navy}QV.13}\quad A loan model never sees religion or caste but splits first on pin code, and approval rates differ sharply by community. What is this called?\hfill\btl{L5}\par\smallskip
{\small \textbf{\color{navyL}(A)}\ Overfitting\\[1pt]\textbf{\color{navyL}(B)}\ Proxy discrimination\\[1pt]\textbf{\color{navyL}(C)}\ Label noise\\[1pt]\textbf{\color{navyL}(D)}\ Concept drift}
\end{tcolorbox}
\begin{tcolorbox}[enhanced,breakable,colback=white,colframe=navyL,arc=2mm,boxrule=0.7pt,left=3mm,right=3mm,top=2mm,bottom=2mm]
\textbf{\color{navy}QV.14}\quad Which expert-system component answers the user's 'WHY are you asking this?' and 'HOW did you conclude that?'\hfill\btl{L2}\par\smallskip
{\small \textbf{\color{navyL}(A)}\ Knowledge acquisition\\[1pt]\textbf{\color{navyL}(B)}\ Inference engine\\[1pt]\textbf{\color{navyL}(C)}\ Explanation facility\\[1pt]\textbf{\color{navyL}(D)}\ Working memory}
\end{tcolorbox}
\begin{tcolorbox}[enhanced,breakable,colback=white,colframe=navyL,arc=2mm,boxrule=0.7pt,left=3mm,right=3mm,top=2mm,bottom=2mm]
\textbf{\color{navy}QV.15}\quad How many leaves does the ID3 tree learned from the 14-example Play Tennis data have?\hfill\btl{L3}\par\smallskip
{\small \textbf{\color{navyL}(A)}\ 3\\[1pt]\textbf{\color{navyL}(B)}\ 4\\[1pt]\textbf{\color{navyL}(C)}\ 5\\[1pt]\textbf{\color{navyL}(D)}\ 7}
\end{tcolorbox}
\newpage
\subsection*{Answer key --- Unit V quiz}
\begin{longtable}{>{\leavevmode\bfseries\color{navy}}p{1.3cm} c p{13.2cm}}
\rowcolor{navy}\textcolor{gold}{\textbf{Q.}}&\textcolor{gold}{\textbf{Ans}}&\textcolor{gold}{\textbf{Worked explanation}}\\\endhead
QV.01&\textbf{\color{leaf}C}&--(9/14) log$_2$(9/14) -- (5/14) log$_2$(5/14) $\approx$ 0.410 + 0.530 = 0.940.\\
\rowcolor{ivory}QV.02&\textbf{\color{leaf}C}&0.940 -- (5/14)(0.971) -- (4/14)(0) -- (5/14)(0.971) $\approx$ 0.247.\\
QV.03&\textbf{\color{leaf}B}&1 -- (9/14)$^2$ -- (5/14)$^2$ = 1 -- 0.413 -- 0.128 $\approx$ 0.459.\\
\rowcolor{ivory}QV.04&\textbf{\color{leaf}B}&w$\cdot$x = --1 < 0, so $\hat y$ = 0 and the error is +1. Update w $\leftarrow$ w + (1)(1, 1, 1) = (0, 1, 1).\\
QV.05&\textbf{\color{leaf}B}&XOR is not linearly separable. Two hidden units (e.g. OR and NAND) combined by an AND output unit compute XOR.\\
\rowcolor{ivory}QV.06&\textbf{\color{leaf}A}&Precision 40/50 = 0.80; recall 40/60 $\approx$ 0.667; F1 = 2PR/(P+R) $\approx$ 0.727.\\
QV.07&\textbf{\color{leaf}B}&A large train--test gap signals high variance; the shallow tree generalises better.\\
\rowcolor{ivory}QV.08&\textbf{\color{leaf}C}&CF = CF$_1$ + CF$_2$(1 -- CF$_1$) = 0.6 + 0.5 $\times$ 0.4 = 0.80.\\
QV.09&\textbf{\color{leaf}B}&AND takes the minimum premise CF (0.7), multiplied by the rule CF: 0.7 $\times$ 0.8 = 0.56.\\
\rowcolor{ivory}QV.10&\textbf{\color{leaf}B}&MYCIN diagnoses goal-first (backward); XCON configures systems data-first (forward).\\
QV.11&\textbf{\color{leaf}B}&Accuracy is 98\% but recall is 0: the accuracy paradox. Imbalanced problems need class-sensitive metrics.\\
\rowcolor{ivory}QV.12&\textbf{\color{leaf}A}&$\sigma$(0) = 1/(1 + e$^0$) = 0.5; tanh(0) = 0; ReLU(--2) = max(0, --2) = 0.\\
QV.13&\textbf{\color{leaf}B}&Pin code correlates with protected attributes, so the model discriminates through a proxy. Audit group metrics and remove or constrain proxies.\\
\rowcolor{ivory}QV.14&\textbf{\color{leaf}C}&The explanation facility traces the rules fired (HOW) and the goal being pursued (WHY).\\
QV.15&\textbf{\color{leaf}C}&Outlook $\to$ Sunny: Humidity (High $\to$ No, Normal $\to$ Yes); Overcast $\to$ Yes; Rain: Wind (Strong $\to$ No, Weak $\to$ Yes). Five leaves.\\
\end{longtable}
